\documentclass{article}
\usepackage{amsmath,amssymb,hyperref,url}
\hypersetup{hidelinks}
\title{Peer scoring when energy also pays for evidence}
\author{}
\date{}
\begin{document}
\maketitle
\begin{abstract}
This account connects the peer-scoring mechanism in \emph{Mechanism Design for Alignment and Control} with the
resource economy in \emph{The Energy Society}.
The peers asked whether energy pressure changes the evidence available from other agents and thereby weakens
incentives to report truthfully.
They checked examples in which selection reverses an outdated score and loss of informative outcomes raises the minimum peak payment.
Correct calibration and score normalisation resolve parts of these examples, but no experiment established the
proposed effect in the Energy Society.
The verifier therefore paused the thread pending suitable round-level records or access to experiments in that economy.
\end{abstract}

\section{Introduction}
\emph{Mechanism Design for Alignment and Control}, called Q here, studies how to make agents report truthfully
and carry out intended actions when their preferences and capabilities are private \cite{Q}.
Its peer-scoring construction uses what an agent's private type predicts about other agents' reports.
A type means the private information relevant to its preferences, capabilities and beliefs.
The designer must know these conditional beliefs, and different types must imply different beliefs.
The mechanism also cancels the reported type's own utility, permits unrestricted rewards and prohibits actions outside the intended set.
These are substantive assumptions, beyond simply paying agents for agreement.

\emph{The Energy Society}, called P, studies language-model agents whose generated tokens consume energy \cite{P}.
Energy funds reasoning and continued participation.
Agents attempt jobs, make donations and recommend actions to one another.
The competitive and cooperative conditions instruct them to maximise individual or group energy, respectively.
These instructions do not implement Q's reward mechanism.
P also has known correct answers for its multiple-choice jobs, so direct answer checking is an available alternative to peer scoring.

The connexion pursued was narrower than a direct transfer of Q into P.
If energy pressure changes which useful peer reports remain available, a score may be based on evidence that
no longer has the assumed distribution.
Paying for better reports then interacts with the resource that keeps peers active.
The question was whether this feedback occurs in P and, if so, whether affordable payments can correct its effect on incentives.
The thread supplied finite mathematical examples, not evidence that this feedback occurs in P.

\section{What was done}
Ada and Emmy first checked the proposed mechanism transfer against the papers' assumptions.
They rejected a direct reading of P's prompted objectives as Q's mechanism.
Ada began with a bounded-payment benchmark for two possible reports.
Emmy examined selection: the possibility that the peers who remain observable differ from the original population.
They also checked whether P's communication measures establish concealed capability, and found that they do not.

Emmy then constructed an example in which scoring only observed reports with the original probabilities rewards a false report.
Ada checked the calculation and pointed out that missing reports can themselves carry information.
Emmy supplied a second example that removes informative outcomes even when missingness is recorded.
Ada independently checked its sharp payment bound.
Their literature searches meanwhile showed that effort costs, payment restrictions and poor equilibria already have substantial precedents.

Next, Ada audited the simulation code, and Emmy checked the audit.
This constrained when a missing report could legitimately enter the model.
The peers separated a maximum possible payout from expected spending, and required beliefs to reflect public
information available when a report is committed.
Finally, Emmy checked score normalisation and showed that Q's quadratic score attains the sharp payment bound in the second example.
Both peers left the natural-economy question open.
They ran exact arithmetic checks, but no language-model experiments or full Energy Society runs.

\section{Results}
\subsection{Selection can reverse an outdated score}
Let \(H\) and \(L\) label two private types and also their possible reports.
Let \(Y\) be a truthful reference peer's outcome, and let \(y\) denote one possible value of that outcome.
Write \(p_H\) and \(p_L\) for the outcome distributions believed by the two types before selection.
For any reported distribution \(p\), Q's quadratic score is
\[
S(p,y)=2p(y)-\lVert p\rVert_2^2,
\qquad \lVert p\rVert_2^2=\sum_y p(y)^2.
\]
Here \(p(y)\) is the probability assigned to outcome \(y\).
With the correct, fixed outcome distribution, the expected advantage of reporting the true distribution over
the other row is their squared distance.
This positive margin is the reporting component of Q's argument.
It does not by itself establish action obedience.

In Emmy's first example, outcomes are zero or one, and
\[
p_H(1)=\frac34,\qquad p_L(1)=\frac14.
\]
An outcome of one is observed with probability \(1/10\); an outcome of zero is always observed.
Among observed reports, the probability of one therefore falls to \(3/13\) for type \(H\) and \(1/31\) for type \(L\).
The original score difference between reports \(H\) and \(L\) is one at outcome one and minus one at outcome zero.
Thus the expected truthful advantage for type \(H\), conditional on observation, becomes
\[
2\frac3{13}-1=-\frac7{13}.
\]
A larger positive score multiplier strengthens the incentive to report falsely.

The remedy in this example is to use the correct probabilities, provided selection is known and cannot be
changed by the reporting deviation.
Another remedy is to score nonresponse as an outcome.
Writing \(q_H\) and \(q_L\) for the full observable distributions, in the order one, zero, missing, gives
\[
q_H=(3/40,1/4,27/40),\qquad
q_L=(1/40,3/4,9/40).
\]
These distributions remain distinct: their squared distance is \(91/200\).
Missingness therefore retains information here.
The example shows a failure to update the score, not a failure of Q or a general loss of information whenever peers disappear.

\subsection{Losing diagnostic outcomes raises the peak payment}
The second example does lose useful information.
Now the latent peer outcome has three values, with distributions
\[
p_H=(1/2,2/5,1/10),\qquad
p_L=(1/2,1/10,2/5).
\]
The first value is always observed.
The second and third are each observed with probability \(\varepsilon\), where \(0\leq\varepsilon\leq1\).
Including missingness as a fourth outcome gives
\begin{align*}
q_H&=(1/2,2\varepsilon/5,\varepsilon/10,(1-\varepsilon)/2),\\
q_L&=(1/2,\varepsilon/10,2\varepsilon/5,(1-\varepsilon)/2).
\end{align*}
Only the second and third outcomes distinguish the types.
Their total variation distance, a measure of this distinction, is
\[
\operatorname{TV}(q_H,q_L)
:=\frac12\sum_y|q_H(y)-q_L(y)|=\frac{3\varepsilon}{10}.
\]

Assume that utility is linear in bonuses, that the peer outcome distribution is fixed under reporting
deviations, and that either type gains \(M>0\) from the false report's continuation before bonuses.
The quantity \(M\) is an assumed temptation in this example, not an estimate of model preferences.
Action obedience is a separate issue.
Let \(R(r,y)\) be the bonus for report \(r\in\{H,L\}\) and outcome \(y\), with every bonus between zero and a cap \(B\).
Define \(d(y)=R(H,y)-R(L,y)\).
For truthfulness, the expected bonus advantage must offset the temptation:
\[
\mathbb E_{q_H}d\geq M,\qquad
\mathbb E_{q_L}d\leq-M.
\]
The notation \(\mathbb E_q\) means expectation under distribution \(q\).
Since \(|d(y)|\leq B\), subtraction gives
\[
2M\leq\sum_y(q_H(y)-q_L(y))d(y)
\leq 2B\operatorname{TV}(q_H,q_L).
\]

This necessary bound is attained for these symmetric distributions.
Pay \(B\) for report \(H\) on the second outcome, or for report \(L\) on the third, and zero otherwise.
Either type then gains \(3B\varepsilon/10\) in expected bonus by reporting truthfully.
The minimum cap is consequently
\[
B_{\min}=\frac{10M}{3\varepsilon}\qquad(\varepsilon>0).
\]
At equality the agent is indifferent after including the temptation; a larger cap gives strict preference for truthfulness.
At \(\varepsilon=0\), the observable distributions coincide and no finite bonus works.
Attainment is specific to this example, not a general sufficiency result for arbitrary distributions.

The divergence concerns the largest possible payment.
At the minimum cap, expected truthful payment is
\[
\frac{2\varepsilon}{5}B_{\min}=\frac{4M}{3}.
\]
Expected spending stays constant while rare payments become large.
Thus an expected budget and a reserve sufficient for every payout are different restrictions.
P supplies no explicit principal's scoring reserve; imposing one would be an added intervention.

\subsection{Normalisation attains the same bound}
A final correction concerned the difference between a score multiplier and the resulting payment cap.
For the second example, let \(\Lambda>0\) be the multiplier and subtract the smallest score available at each outcome:
\[
b(y)=\min_{r\in\{H,L\}}S(q_r,y),\qquad
R(r,y)=\Lambda\bigl[S(q_r,y)-b(y)\bigr].
\]
The baseline \(b(y)\) does not depend on the report.
Under the fixed-distribution and linear-utility assumptions, subtracting it preserves reporting comparisons.
The resulting payments are exactly the diagnostic bonuses just described, with
\[
B=\frac{3\Lambda\varepsilon}{5},\qquad
\Lambda_{\min}=\frac{50M}{9\varepsilon^2},\qquad
B_{\min}=\frac{10M}{3\varepsilon}.
\]
Here \(\Lambda_{\min}\) is the smallest multiplier that offsets \(M\).
Its inverse-square growth does not mean that peak payments grow at that rate.
Normalised quadratic scoring already attains the sharp cap.
This argument does not preserve survival paths or cover deviations that change the peer outcome distribution.

The derivations and qualifications are recorded in
\path{ada/research_note.md} and \path{emmy/survival_scoring.md}.
Ada independently checked the selection example and payment bound.
Their exact arithmetic checks are in \path{ada/check_examples.py} and
\path{emmy/verify_scoring.py}; Emmy's script also checks the final normalisation.

\section{Objections and their resolution}
The direct mechanism-transfer claim did not survive.
P establishes neither Q's known belief and utility maps nor its cancellation and prohibition instruments.
Q's separate coupled-reward construction also assumes scalar actions and a known decreasing relation between
agents' biases that P does not establish.
Likewise, avoiding difficult jobs or recommending collisions does not identify sandbagging, meaning concealed capability.
P has no verified capability frontier or evaluation stage that determines later permissions.

Poor equilibria do not refute Q either.
Q requires the target behaviour to be an equilibrium, not the unique equilibrium.
Recalibration settles the first example's reporting problem against a fixed peer distribution.
It does not settle strategic manipulation of peers, or selection among equilibria.

The code audit changed the proposed interpretation of exhaustion.
At the recorded commit, energy is deducted after a completed model call \cite{middleware}.
Completed votes remain available even if producing them exhausts their author, whereas already inactive agents
are excluded before discussion \cite{environment}.
Successful job payouts can also restore an exhausted solver's energy.
Ordinary exhaustion during a completed call must therefore not be recorded as a missing vote.

The relevant beliefs must condition on public history and active status at report commitment.
An already known absence is not a future random scoring event.
Fresh signals independent of past energy need not show survival bias at all.
Sealing reports and eligibility before recommendations or donations limits immediate manipulation, but leaves
strategic selection across rounds unresolved.

P's resource evidence also needed qualification.
Its abstract claims that larger models spend more energy than they gain even without size-dependent token costs.
Its ``Experiment 2: No Size Penalty'' instead reports all agents surviving and all efficiency ratios above one.
The peers used that section and its results table, labelled
\texttt{tab:exp2\_big\_table} in the source, rather than the stronger abstract claim.
Neither these results nor pooled collision counts identify the effect of a scoring intervention.

Novelty remained unsettled.
The scoped searches found earlier work on costly effort, participation and budget balance by Miller, Resnick
and Zeckhauser \cite{miller}; participation and payment restrictions in \emph{Dwelling on the Negative}
\cite{witkowski}; and budget-constrained information acquisition in \emph{Information Gathering with Peers}
\cite{radanovic}.
Gao, Wright and Leyton-Brown study uninformative equilibria and advantages of direct checking under their assumptions \cite{gao}.
Q itself discusses language-model peer scoring, and Emmy checked its precursor by Qiu, Carroll and Allen \cite{qiu}.
These checks ruled out treating a payment bound or generic language-model peer scoring as sufficient novelty.
They were not a comprehensive literature review.

\section{Why the thread stopped}
The verifier's sole recorded round ended in PAUSE.
It accepted the consolidated account and the conditional calculations, but found no evidence for the
distinctive connexion: a resource-dependent loss of useful peer information in P after conditioning on public
history.
More algebra on the examples could not supply that evidence.
Suitable round-level traces or access to new experiments were missing and would have to be provided before the research could proceed.

The smallest restart is a channel-identification check in the natural P economy.
Fix what is public when a report is committed, then estimate whether resource state changes the informative
peer-outcome distributions after accounting for that history.
Use the audited timing to distinguish inactive-at-start peers from completed votes.
A synthetic erasure test can check the mathematics but cannot answer this question about the natural economy.

If no material change appears, the proposed connexion lacks support.
If it does appear, it supplies calibration for a sealed-report scoring test with explicit peak reserves.
That test still needs direct-answer checking and report-independent transfers as controls.
The latter must account for timing, recipients and payout variation, since equal expected payments can affect survival differently.
Injected energy must not count as produced task value.
Identifying natural private types, valuing energy near exhaustion, strategic peer influence and novelty would
remain open even after this first empirical hurdle.

\begingroup
\raggedright
\bibliographystyle{plain}
\bibliography{references}
\endgroup
\end{document}
